\section*{Hoofdstuk \ref{ch:b2:unicellular-diversity} --- Diversiteit van eencellige organismen}

\begin{solution}{exo:b2:unicellular-diversity:1}
\omterm{def:b2:unicellular-diversity:unicellular}{Eencellig}: één cel voert alle functies uit en leeft alleen; \omterm{def:b2:unicellular-diversity:unicellular}{koloniaal}:
identieke cellen blijven aan elkaar zitten, maar elk zou alleen kunnen
leven; \omterm{def:b2:unicellular-diversity:unicellular}{meercellig}: verschillende celtypen, onderling afhankelijk, alleen
de kiembaan plant zich voort. \emph{E.~coli} en de gist zijn \omterm{def:b2:unicellular-diversity:unicellular}{eencellig}
(een knop laat los en leeft alleen); \emph{Gonium} is \omterm{def:b2:unicellular-diversity:unicellular}{koloniaal} (16
identieke cellen, elk in staat een kolonie te stichten); \emph{Volvox}
is \omterm{def:b2:unicellular-diversity:unicellular}{meercellig} (sterfelijke somatische cellen, enkele kiemcellen).
\end{solution}

\begin{solution}{exo:b2:unicellular-diversity:2}
$S/V = 3/r$: \qty{6}{\micro m^{-1}} voor de coccus,
\qty{0.06}{\micro m^{-1}} voor de \omterm{def:b2:unicellular-diversity:protist}{protist}, een verschil van een factor honderd.
De coccus ademt sneller per volume-eenheid: elk deel van zijn
cytoplasma ligt binnen een fractie van een micrometer van het oppervlak
waardoor zuurstof en voedsel binnenkomen, terwijl het inwendige van de
\omterm{def:b2:unicellular-diversity:protist}{protist} per volume-eenheid via honderd keer minder oppervlak wordt gevoed.
\end{solution}

\begin{solution}{exo:b2:unicellular-diversity:3}
Geen kern (een \omterm{def:b2:unicellular-diversity:prokaryote}{nucleoïde}) tegenover een kernenvelop; geen mitochondriën of
endomembranen, met de ademhalingsketen in het plasmamembraan, tegenover
organellen; een flagel die een draaiende stijve schroef van flagelline is,
door protonen aangedreven, tegenover een trilhaar met $9+2$-structuur dat
buigt door dyneïne en ATP; verder de grootte (\qty{1}{\micro m} tegenover
tientallen) en een wand van peptidoglycaan. Een bacterie kan niet
fagocyteren: haar wand verhindert het omsluiten van een deeltje, wat een
amoebe met haar schijnvoetjes wel doet.
\end{solution}

\begin{solution}{exo:b2:unicellular-diversity:4}
Trilharen: voortbeweging en het toewuiven van voedsel naar de mond;
mondgroef: leidt deeltjes naar de celmond; voedselvacuole: verteert
opgenomen bacteriën met lysosomale enzymen; \omterm{prop:b2:unicellular-diversity:osmoregulation}{contractiele vacuole}: verdrijft
het water dat door osmose binnenkomt; macronucleus: de werkende kern,
honderden genkopieën die voor het dagelijkse leven worden overgeschreven;
micronucleus: de generatieve diploïde kern, gebruikt bij meiose en conjugatie.
\end{solution}

\begin{solution}{exo:b2:unicellular-diversity:5}
$t = x^2/2D$: $(10^{-6})^2/10^{-9} = \qty{1}{ms}$ voor de bacterie;
$(5\times 10^{-4})^2/10^{-9} = \qty{250}{s}$, ongeveer vier minuten,
voor de amoebe. De amoebe kan niet op diffusie vertrouwen om metabolieten
te verdelen: ze heeft cytoplasmastroming en door motoreiwitten aangedreven
transport langs het cytoskelet nodig, waar de bacterie buiten kan.
\end{solution}

\begin{solution}{exo:b2:unicellular-diversity:6}
$Re = \rho vL/\eta$: \emph{Paramecium} $10^3\times 10^{-3}\times
2\times 10^{-4}/10^{-3} = 0.2$; bacterie $10^3\times 2\times
10^{-5}\times 2\times 10^{-6}/10^{-3} = 4\times 10^{-5}$; zwemmer
$10^3\times 1\times 2/10^{-3} = 2\times 10^6$. Bij $Re \gg 1$ draagt
de traagheid de zwemmer na de slag verder; bij $Re \ll 1$ stopt de
viskeuze wrijving de bacterie binnen een nanometer, zodat ze alleen
beweegt zolang ze slaat.
\end{solution}

\begin{solution}{exo:b2:unicellular-diversity:7}
$v_s = 2r^2\Delta\rho g/9\eta = 2\times 10^{-10}\times 100\times
9.81/(9\times 10^{-3}) = \qty{2.2e-5}{m/s}$; \qty{50}{m} kost
$50/2.2\times 10^{-5} = \qty{2.3e6}{s}$, 27 dagen. Halveren van de straal
deelt $v_s$ door 4: 108 dagen. Verlagen van $\Delta\rho$ tot
\qty{20}{kg/m^3} deelt haar door 5: 135 dagen.
\end{solution}

\begin{solution}{exo:b2:unicellular-diversity:8}
$V = \tfrac{4}{3}\pi\times 100\times 25\times 25 =
\qty{2.6e5}{\micro m^3}$. In \qty{1800}{s}: $2.6\times 10^5/1800 =
\qty{145}{\micro m^3/s}$; omdat $\qty{1}{\micro m^3} =
\qty{1e-15}{L}$, is dit \qty{1.5e-13}{L/s}.
\end{solution}

\begin{solution}{exo:b2:unicellular-diversity:9}
Een verdubbeling over \qty{1}{mm} is een stijging van ongeveer
\qty{0.07}{\%} per micrometer (want $2^{1/1000} - 1 = 0.0007$), dus
\qty{0.07}{\%} over de cel. Een run van \qty{1}{s} legt \qty{20}{\micro m} af:
\qty{1.4}{\%}. Om een verschil van \qty{0.07}{\%} te onderscheiden, zou de
cel aan elk uiteinde zo'n $(1/0.0007)^2 = 2\times 10^6$ moleculen moeten
tellen, veel meer dan er in een seconde aan haar paar duizend receptoren
binden; \qty{1.4}{\%} vraagt er ongeveer \num{5000}, wat haalbaar is. De
vergelijking in de tijd maakt van een hopeloze ruimtelijke meting een
mogelijke, doordat de run het integreren doet.
\end{solution}

\begin{solution}{exo:b2:unicellular-diversity:10}
Oppervlak $4\pi r^2$ met $r = \qty{250}{\micro m}$: \qty{7.9e5}{\micro
m^2}; cytoplasmatisch volume $\approx$ oppervlak $\times$ \qty{2}{\micro m}
$= \qty{1.6e6}{\micro m^3}$; verhouding \qty{0.5}{\micro m^{-1}}. Voor
\emph{E.~coli}, een cilinder: $S/V = 2\pi r L/\pi r^2 L = 2/r =
\qty{4}{\micro m^{-1}}$. De reus doet het maar acht keer slechter dan
\emph{E.~coli}, niet 500 keer, omdat elk punt van zijn cytoplasma
binnen \qty{2}{\micro m} van het oppervlak ligt. De vacuole slaat
nitraat op, zijn elektronenacceptor, zodat hij sulfide kan verademen in
het zuurstofloze sediment tussen de zeldzame keren dat het water wordt
omgeroerd en nitraat aanvoert (\cref{ch:b2:microbial-metabolism}).
\end{solution}

\begin{solution}{exo:b2:unicellular-diversity:11}
Een cel kan niet tegelijk zwemmen en zich delen. Een kolonie van 16
cellen die tijdens de uren van de deling stopt met zwemmen, zakt volgens
de \omterm{thm:b2:unicellular-diversity:stokes}{wet van Stokes} een verwaarloosbare afstand (haar straal is enkele
tientallen micrometers, haar zinksnelheid micrometers per seconde); de
hele kolonie kan zich delen zonder iets te verliezen. Een bol van
\num{10000} cellen met straal \qty{500}{\micro m} zou honderd keer
sneller zinken --- millimeters per seconde, meters per uur --- en uit
het licht verdwijnen: het loont om de meeste cellen te laten zwemmen
terwijl enkele zich delen. De kiemcellen zijn groot omdat ze de
middelen moeten meedragen om zonder te eten, door opeenvolgende delingen,
een hele dochterbol van duizenden cellen op te bouwen; ze worden dus
door het soma bevoorraad.
\end{solution}

\begin{solution}{exo:b2:unicellular-diversity:12}
Een clade is een voorouder met al zijn nakomelingen; een graad is een
organisatieniveau dat door verscheidene lijnen is bereikt. \omterm{def:b2:unicellular-diversity:protist}{Protisten} ---
groene algen (bij de landplanten), kiezelwieren (bij de bruinwieren),
ciliaten (bij de dinoflagellaten en \emph{Plasmodium}), amoeben (dicht
bij dieren en schimmels) --- delen alleen de eukaryote cel die alleen
leeft, en hun laatste gemeenschappelijke voorouder is die van alle
eukaryoten, ons inbegrepen: een graad. ``\omterm{def:b2:unicellular-diversity:prokaryote}{Prokaryoot}'' verenigt bacteriën
en archaea op grond van wat ze missen; de archaea staan dichter bij de
eukaryoten dan bij de bacteriën, dus is ook deze term een graad.
Cyanobacteriën stammen af van één voorouder die de zuurstofproducerende
fotosynthese uitvond en omvatten al zijn nakomelingen (afgezien van de
lijn van de chloroplast): een clade.
\end{solution}

\begin{solution}{pb:b2:unicellular-diversity:1}
\textbf{1.} $1\times 1\times 0.1 = \qty{0.1}{mm^3} = \qty{0.1}{\micro L}$.
\textbf{2.} $24/\qty{0.1}{\micro L} = 240$ per \unit{\micro L} $=
\qty{2.4e5}{}$ cellen per milliliter.
\textbf{3.} $V = \tfrac{4}{3}\pi\times 125 = \qty{524}{\micro m^3}$;
$S/V = 3/5 = \qty{0.6}{\micro m^{-1}}$.
\textbf{4.} $2.4\times 10^5\times 524 = \qty{1.26e8}{\micro m^3}$ per
milliliter, en $\qty{1}{mL} = \qty{1e12}{\micro m^3}$: een fractie
$1.3\times 10^{-4}$, \qty{0.013}{\%}.
\textbf{5.} $\qty{1000}{m^3} = \qty{1e9}{mL}$: $2.4\times 10^{14}$
cellen.
\textbf{6.} De cyanobacterie is een draad van veel kleinere cellen
(enkele micrometers), blauwgroen en zonder duidelijke chloroplast of kern,
en ze zwemt niet met flagellen; de alg is één eivormige cel met twee
flagellen, een oogvlek en één bekervormige chloroplast.
\textbf{7.} $384/24 = 16 = 2^4$: vier verdubbelingen in \qty{48}{h},
$T = \qty{12}{h}$.
\textbf{8.} $k = \ln 2/T = 0.693/12 = \qty{0.058}{h^{-1}}$.
\textbf{9.} $t = \ln(10^8/2.4\times 10^5)/k = \ln(417)/0.058 =
6.03/0.058 = \qty{104}{h}$, ongeveer 4.3 dagen.
\textbf{10.} Voedingsstoffen (fosfaat, nitraat) raken op; de cellen
beschaduwen elkaar, zodat het licht per cel afneemt; grazers (ciliaten,
raderdiertjes, \emph{Daphnia}) vermenigvuldigen zich op hun kosten; bij
\qty{1e8}{} per milliliter zouden de cellen zo'n vijf procent van het volume
vullen en het opgeloste $\mathrm{CO_2}$ uitputten.
\textbf{11.} Groei gedurende slechts de helft van de tijd: de gemiddelde
snelheid wordt gehalveerd, \qty{208}{h}, 8.7 dagen.
\textbf{12.} De telling meet de populatie, niet de afzonderlijke cellen:
als een cel om de \qty{36}{h} 8 dochtercellen vrijlaat, wordt de populatie
met 8 $= 2^3$ vermenigvuldigd in drie \omterm{def:b2:unicellular-diversity:division}{generatietijden} van elk \qty{12}{h}.
De \omterm{def:b2:unicellular-diversity:division}{generatietijd} is de gemiddelde verdubbelingstijd van het aantal, hoe
de delingen ook gegroepeerd zijn.
\textbf{13.} $Re = 10^3\times 10^{-4}\times 10^{-5}/10^{-3} = 10^{-3}$.
\textbf{14.} $m = 1050\times 5.24\times 10^{-16} = \qty{5.5e-13}{kg}$;
$6\pi\eta r = 6\pi\times 10^{-3}\times 5\times 10^{-6} =
\qty{9.4e-8}{kg/s}$; $d = 5.5\times 10^{-13}\times 10^{-4}/9.4\times
10^{-8} = \qty{5.8e-10}{m}$: een halve nanometer.
\textbf{15.} $v_s = 2\times 25\times 10^{-12}\times 50\times
9.81/(9\times 10^{-3}) = \qty{2.7e-6}{m/s}$, \qty{2.7}{\micro m/s}.
\textbf{16.} \qty{1}{m} zinken: $1/2.7\times 10^{-6} =
\qty{3.7e5}{s}$, 4.3 dagen; omhoog zwemmen met \qty{100}{\micro m/s}:
\qty{1e4}{s}, 2.8 uur. Zwemmen wint met een factor 37.
\textbf{17.} $c_\infty = \qty{1e-2}{mol/m^3}$: $J = 4\pi\times
1.9\times 10^{-9}\times 5\times 10^{-6}\times 10^{-2} =
\qty{1.2e-15}{mol/s} = 7.2\times 10^8$ moleculen per seconde.
\textbf{18.} \qty{100}{pg} $= \qty{1e-10}{g} = \qty{8.3e-12}{mol}$
koolstof, verdubbeld in \qty{12}{h} $= \qty{43200}{s}$:
\qty{1.9e-16}{mol/s} $= 1.2\times 10^8$ atomen per seconde. Aanvoer
gedeeld door vraag: $7.2/1.2 = 6$.
\textbf{19.} De aanvoer wordt tien keer zo groot ($\propto r$), de vraag
duizend keer ($\propto r^3$): de verhouding wordt $6/100 = 0.06$; zo'n
cel zou maar met zes procent van de snelheid kunnen groeien, dus een door
diffusie gevoede alg van die grootte is onmogelijk zonder omroeren,
inwendig transport of een veel trager metabolisme.
\textbf{20.} $\Delta\Pi = RT\Delta c = 8.314\times 293\times 99 =
\qty{2.4e5}{Pa}$, \qty{2.4}{bar}.
\textbf{21.} Oppervlak $4\pi(5\times 10^{-6})^2 = \qty{3.14e-10}{m^2}$;
instroom $L_p A\Delta\Pi = 10^{-14}\times 3.14\times 10^{-10}\times
2.4\times 10^5 = \qty{7.6e-19}{m^3/s} = \qty{0.76}{\micro m^3/s}$.
\textbf{22.} $524/0.76 = \qty{690}{s}$, ongeveer elf minuten.
\textbf{23.} $P = \Delta\Pi\times$ debiet $= 2.4\times 10^5\times
7.6\times 10^{-19} = \qty{1.8e-13}{W}$.
\textbf{24.} Eén ATP: $5\times 10^4/6.02\times 10^{23} =
\qty{8.3e-20}{J}$: $1.8\times 10^{-13}/8.3\times 10^{-20} = 2.2\times
10^6$ ATP per seconde. Koolstoffixatie: $3\times 1.2\times 10^8 =
3.6\times 10^8$ ATP per seconde. De vacuole kost \qty{0.6}{\%} daarvan
--- een goedkope verzekering tegen barsten.
\textbf{25.} Waterinstroom \qty{0.76}{\micro m^3/s} (het volume van de
cel in elf minuten); zonder \omterm{prop:b2:unicellular-diversity:osmoregulation}{contractiele vacuole} zou de cel haar volume in
ongeveer \qty{700}{s} verdubbelen; droog blijven kost minstens
$2\times 10^6$ ATP per seconde, minder dan een procent van wat haar
fotosynthese aan koolstof besteedt.
\end{solution}
